\section{Reproducibility Checklist for JAIR}

Select the answers that apply to your research -- one per item.

\subsection*{All articles:}

\begin{enumerate}
    \item All claims investigated in this work are clearly stated.
    [yes]
    \item Clear explanations are given how the work reported substantiates the claims.
    [yes]
    \item Limitations or technical assumptions are stated clearly and explicitly.
    [yes]
    \item Conceptual outlines and/or pseudo-code descriptions of the AI methods introduced in this work are provided, and important implementation details are discussed.
    [yes]
    \item
    Motivation is provided for all design choices, including algorithms, implementation choices, parameters, data sets and experimental protocols beyond metrics.
    [yes]
\end{enumerate}

\subsection*{Articles containing theoretical contributions:}
Does this paper make theoretical contributions?
[no]

If yes, please complete the list below.

\begin{enumerate}
    \item All assumptions and restrictions are stated clearly and formally.
    [yes/partially/no]
    \item All novel claims are stated formally (e.g., in theorem statements).
    [yes/partially/no]
    \item Proofs of all non-trivial claims are provided in sufficient detail to permit verification by readers with a reasonable degree of expertise (e.g., that expected from a PhD candidate in the same area of AI). [yes/partially/no]
    \item
    Complex formalism, such as definitions or proofs, is motivated and explained clearly.
    [yes/partially/no]
    \item
    The use of mathematical notation and formalism serves the purpose of enhancing clarity and precision; gratuitous use of mathematical formalism (i.e., use that does not enhance clarity or precision) is avoided.
    [yes/partially/no]
    \item
    Appropriate citations are given for all non-trivial theoretical tools and techniques.
    [yes/partially/no]
\end{enumerate}

\subsection*{Articles reporting on computational experiments:}
Does this paper include computational experiments? [yes]

If yes, please complete the list below.
\begin{enumerate}
    \item
    All source code required for conducting experiments is included in an online appendix
    or will be made publicly available upon publication of the paper.
    The online appendix follows best practices for source code readability and documentation as well as for long-term accessibility.
    [yes (public GitHub repository)]
    \item The source code comes with a license that
    allows free usage for reproducibility purposes.
    [no]
    \item The source code comes with a license that
    allows free usage for research purposes in general.
    [no]
    \item
    Raw, unaggregated data from all experiments is included in an online appendix
    or will be made publicly available upon publication of the paper.
    The online appendix follows best practices for long-term accessibility.
    [yes]
    \item The unaggregated data comes with a license that
    allows free usage for reproducibility purposes.
    [no]
    \item The unaggregated data comes with a license that
    allows free usage for research purposes in general.
    [no]
    \item If an algorithm depends on randomness, then the method used for generating random numbers and for setting seeds is described in a way sufficient to allow replication of results.
    [yes]
    \item The execution environment for experiments, the computing infrastructure (hardware and software) used for running them, is described, including GPU/CPU makes and models; amount of memory (cache and RAM); make and version of operating system; names and versions of relevant software libraries and frameworks.
    [partially]
    \item
    The evaluation metrics used in experiments are clearly explained and their choice is explicitly motivated.
    [yes]
    \item
    The number of algorithm runs used to compute each result is reported.
    [yes]
    \item
    Reported results have not been ``cherry-picked'' by silently ignoring unsuccessful or unsatisfactory experiments.
    [yes]
    \item
    Analysis of results goes beyond single-dimensional summaries of performance (e.g., average, median) to include measures of variation, confidence, or other distributional information.
    [yes]
    \item
    All (hyper-) parameter settings for
    the algorithms/methods used in experiments have been reported, along with the rationale or method for determining them.
    [yes]
    \item
    The number and range of (hyper-) parameter settings explored prior to conducting final experiments have been indicated, along with the effort spent on (hyper-) parameter optimisation.
    [NA]
    \item
    Appropriately chosen statistical hypothesis tests are used to establish statistical significance
    in the presence of noise effects.
    [yes]
\end{enumerate}


\subsection*{Articles using data sets:}
Does this work rely on one or more data sets (possibly obtained from a benchmark generator or similar software artifact)?
[yes]

If yes, please complete the list below.
\begin{enumerate}
    \item
    All newly introduced data sets
    are included in an online appendix
    or will be made publicly available upon publication of the paper.
    The online appendix follows best practices for long-term accessibility with a license
    that allows free usage for research purposes.
    [yes]
    \item The newly introduced data set comes with a license that
    allows free usage for reproducibility purposes.
    [no]
    \item The newly introduced data set comes with a license that
    allows free usage for research purposes in general.
    [no]
    \item All data sets drawn from the literature or other public sources (potentially including authors' own previously published work) are accompanied by appropriate citations.
    [NA]
    \item All data sets drawn from the existing literature (potentially including authors’ own previously published work) are publicly available. [NA]
    \item All new data sets and data sets that are not publicly available are described in detail, including relevant statistics, the data collection process and annotation process if relevant.
    [yes]
    \item
    All methods used for preprocessing, augmenting, batching or splitting data sets (e.g., in the context of hold-out or cross-validation)
    are described in detail. [yes]
\end{enumerate}

\subsection*{Explanations on any of the answers above (optional):}

The study introduces an evaluation protocol rather than a new learning algorithm; its procedures are given in Section~3, Figure~1 and Appendix~B. Decoding settings followed each model card and were fixed in the preregistration; no hyperparameters were tuned (computational item 14). The new data sets are the model generations, judge codes and anonymised human codes released in the repository; no external data sets were used (data items 4--5). All preregistered and exploratory analyses are reported, including null results and deviations (Appendix~A). Licence checklist items are answered ``no'' because repository licence files were not added in this submission build (Phase~20 §0/§5: licence release not approved at build time).
